% Replace appendix/experimental_cost.tex. All existing reference labels
% and numerical results are retained. No new profiling is claimed.
\section{Computational cost and timing boundaries}
\label{app:parameter_counts}

The computational cost of a sparse reduced model depends on both the
parameters selected by its routers and the implementation used to
advance and reconstruct the flow. We therefore examine parameter
utilization and measured prediction time separately. The Circular H/P
specialists provide the parameter comparison, while Circular Periodic
is used to compare local prediction costs and to place the measured
ROM time in the context of a historical full-order simulation.

\subsection{Parameter utilization}

Table~\ref{tab:parameter_counts} compares the total trainable capacity
of each specialist with the parameters active in one prediction-path
routing decision. The active count includes the encoder and refinement
blocks, executed routers, the selected group's shared experts,
channel-specific Top-2 experts, and learned pressure correction.
Each trainable parameter is counted once. Fixed POD bases, projected
operators, normalization statistics, and non-trainable buffers are
excluded. The count refers to a batch-size-one decision, rather than
the union of parameters used over a trajectory or all integration stages.

\begin{table}[tbp]
\centering
\caption{Total trainable capacity and per-decision active parameters
of the Circular specialists. The active fraction measures parameter
utilization, not the fraction of total runtime or memory consumption.}
\label{tab:parameter_counts}
\begin{tabular}{@{}lrrr@{}}
\toprule
Specialist & Total trainable & Active prediction & Active fraction\\
\midrule
Hopf & 31,810,392 & 14,165,816 & 44.53\%\\
Periodic & 48,617,547 & 8,308,959 & 17.09\%\\
\bottomrule
\end{tabular}
\end{table}

The Periodic specialist has more total trainable parameters than Hopf
but uses fewer in each analyzed decision. This illustrates the
distinction between available model capacity and the capacity selected
for an individual prediction. It does not establish a corresponding
reduction in floating-point operations or wall-clock time. In particular,
the Dense control is matched to the active neural parameter budget,
not to the full sparse model's parameter count. A smaller active fraction
therefore does not, by itself, imply less computation than this control.

\subsection{Local prediction time}
\label{app:cost_local_latency}

To test whether sparse parameter use translates into lower latency,
we compare the native Circular Periodic implementation with Dense and
Structured for seed 1248, batch size one, and a $K=48$ rollout.
Table~\ref{tab:local_latency} reports median times from ten repetitions
following three warm-up runs. The timed prediction path includes modal
advancement, feature generation, data transfers, algebraic pressure
reconstruction, decoding of all predicted fields, and pressure gauge
correction. Physical-history encoding, model and data loading, error
evaluation, and the outer E2/descriptor/T2-C pipeline are excluded.
The measurements thus compare local prediction paths rather than
single network evaluations or the complete hierarchy.

\begin{table}[tbp]
\centering\small
\caption{Circular Periodic local prediction cost at $K=48$.
Times are medians over ten repetitions after three warm-ups.
Memory is peak allocated GPU memory in bytes, excluding host memory;
the last column normalizes latency by Dense.}
\label{tab:local_latency}
\begin{tabular}{@{}lrrr@{}}
\toprule
Model & Median time (s) & GPU peak (bytes) & Time / Dense\\
\midrule
Proposed & 2.1848 & 204,238,336 & 3.32\\
Dense & 0.6579 & 43,111,936 & 1.00\\
Structured & 0.8652 & 43,316,736 & 1.32\\
\bottomrule
\end{tabular}
\end{table}

The sparse implementation takes $3.32$ times as long as Dense and also
has higher peak allocated GPU memory; Structured is intermediate in
latency. Thus the reduced active fraction in
Table~\ref{tab:parameter_counts} does not translate into a measured
runtime or memory advantage over these controls. Parameter matching
does not match the number of executed operations, their scheduling,
or data movement. Although the native Galerkin NumPy path runs on the
CPU, the aggregate timings do not isolate its contribution to the
difference between models. Routing overhead, expert execution, and
CPU/GPU transfers remain possible explanations requiring component-level
profiling, rather than established causes. The result should therefore
be read as the cost of the current implementation, not an inherent
speed ordering of sparse and dense architectures.

\subsection{Comparison with a historical full-order simulation}
\label{app:cost_fom_reference}

The comparison with Dense addresses the cost of the local architecture;
a separate comparison with CFD addresses the cost of replacing a
full-order simulation by reduced prediction. For this purpose, we time
the original Circular Periodic specialist with physical-history encoding
and history right-hand-side construction included. At
$\mathrm{Re}=70.314635$, the first fixed valid $K=48$ window extends
from $t=1384.49597$ to $1983.19690$, covering $598.70093$ physical time
units. The input consists of three POD-reconstructed history frames on
97,368 points, and all 48 velocity and pressure predictions are decoded
with pressure gauge correction. The model and inputs are already
resident in memory, so history acquisition and loading are not timed.

On an RTX 4090 with a Xeon Platinum 8358P host, one CPU numerical
thread, and batch size one, ten CUDA-synchronized repetitions after
three warm-ups give a median of $2.18747$ s (mean $2.18729$ s, sample
standard deviation $0.00462$ s, range $2.18034$--$2.19550$ s).
This differs from Table~\ref{tab:local_latency} because the timing
boundary additionally includes the physical-input processing.
The corresponding production OpenFOAM interval contains 4,692 solver
steps and a ClockTime difference of 1,996 s. Its endpoints agree with
the ROM interval to within $2.3\times10^{-5}$ physical time units,
consistent with stored timestamp precision.

The ratio $1996/2.18747\simeq912.5$ shows that this warm local ROM
prediction takes substantially less wall time than the corresponding
historical CFD interval. It is nevertheless not a hardware-matched
benchmark. The FOM ran as one process on a Core i5-14600K VMware host
with 16 vCPUs and at most five concurrent cases; that configuration
was identified retrospectively, and only one FOM timing realization
is available. The FOM interval includes native field writes, whereas
the ROM measurement excludes prediction writes, E2 selection,
descriptor construction, and T2-C fusion. FOM meshing, initialization,
periodic detection, and subsequent VTK/NPZ conversion are outside the
recorded interval, while offline ROM construction and training are
outside the prediction timing.

These differences limit the ratio to a historical FOM/warm-specialist
reference for this case. It neither measures the cost of advancing and
combining two specialists nor establishes a hardware-independent
end-to-end speedup. Taken together, the measurements distinguish two
findings: the local ROM is inexpensive relative to this historical CFD
run, but its current sparse implementation is more expensive than the
Dense and Structured controls. Neither finding implies a runtime result
for Square or Pinball.
